\subsection{Implementasi \textit{From Scratch}}
\label{subsec:dtl-implementasi}

\subsubsection{Landasan Teori}
\todo{Jelaskan entropi, \textit{information gain}, \textit{split information}, \textit{gain ratio}; pencarian \textit{threshold} untuk atribut kontinu; penanganan \textit{missing value} dengan bobot fraksional; dan \textit{pruning} (\textit{pessimistic error pruning} C4.5). Acu \cite{slide-dtl,slide-dtl-issues,quinlan1993c45}. Jelaskan juga alasan memilih C4.5 dibanding ID3 dan CART. Sertakan rumus utama dalam lingkungan \texttt{equation}/\texttt{align}.}

\subsubsection{Struktur Kelas dan Fungsi}
\todo{Deskripsikan kelas pada \texttt{src/algoritma/dtl/}: tiap modul (\texttt{criteria.py}, \texttt{splitting.py}, \texttt{tree.py}, \texttt{pruning.py}, \texttt{serialization.py}, \texttt{dtl.py}) beserta fungsi pentingnya dan parameter kelas \texttt{DTL}. Isi tabel di bawah, satu baris per fungsi/metode.}

\begin{table}[H]
	\centering
	\caption{Deskripsi kelas dan fungsi DTL}
	\label{tab:dtl-fungsi}
	\begin{tabularx}{\textwidth}{l Y}
		\toprule
		\textbf{Fungsi/Metode} & \textbf{Deskripsi} \\
		\midrule
		\texttt{nama\_fungsi} & \todo{deskripsi} \\
		\bottomrule
	\end{tabularx}
\end{table}

\subsubsection{\textit{Source Code}}
\todo{Sisipkan potongan kode inti (perhitungan \textit{gain ratio}, pembangunan pohon rekursif, dan \textit{pruning}) dengan makro \texttt{\textbackslash kode}, lalu jelaskan alurnya. Cek ulang nomor baris jika kode berubah.}
% \kode{../../src/algoritma/dtl/dtl.py}{1}{30}{Judul potongan kode}{lst:dtl-inti}

\subsubsection{\textit{Save} dan \textit{Load} Model}
\todo{Format penyimpanan (JSON/\texttt{.npz}/\textit{pickle}), isi yang disimpan, dan contoh pemakaian.}
